\documentclass[aspectratio=169]{beamer}

%\documentclass[handout]{beamer}
%% To make 4 per page
%\usepackage{pgfpages}
%\mode<handout>{\setbeamercolor{background canvas}{bg=white}}
%\pgfpagesuselayout{4 on 1}[letterpaper,landscape]%,border shrink=5mm]

\usetheme{default}
% Slide setup, colour independent

\usepackage{amsmath,amssymb,amsthm}
\usepackage{colortbl}
\usepackage{bm}
\usepackage{xcolor}
\usepackage{dsfont}
\usepackage{setspace}
%\usepackage{subfigure}
% To use \ding{234} and the like
\usepackage{pifont}
% To cross reference between slide files
\usepackage{zref-xr,zref-user}
% Use something like
% \zexternaldocument{fileI}
% in the tex files. And cite using \zref instead of \ref

% Fields and the like
\def\IC{\mathbb{C}}
\def\IF{\mathbb{F}}
\def\II{\mathbb{I}}
\def\IJ{\mathbb{J}}
\def\IM{\mathbb{M}}
\def\IN{\mathbb{N}}
\def\IP{\mathbb{P}}
\def\IR{\mathbb{R}}
\def\IZ{\mathbb{Z}}
\def\11{\mathds{1}}


% Bold lowercase
\def\ba{\mathbf{a}}
\def\bb{\mathbf{b}}
\def\bc{\mathbf{c}}
\def\bd{\mathbf{d}}
\def\be{\mathbf{e}}
\def\bf{\mathbf{f}}
\def\bh{\mathbf{h}}
\def\bi{\mathbf{i}}
\def\bj{\mathbf{j}}
\def\bk{\mathbf{k}}
\def\bn{\mathbf{n}}
\def\bp{\mathbf{p}}
\def\br{\mathbf{r}}
\def\bs{\mathbf{s}}
\def\bu{\mathbf{u}}
\def\bv{\mathbf{v}}
\def\bw{\mathbf{w}}
\def\bx{\mathbf{x}}
\def\by{\mathbf{y}}
\def\bz{\mathbf{z}}

% Bold capitals
\def\bB{\mathbf{B}}
\def\bD{\mathbf{D}}
\def\bF{\mathbf{F}}
\def\bG{\mathbf{G}}
\def\bI{\mathbf{I}}
\def\bL{\mathbf{L}}
\def\bN{\mathbf{N}}
\def\bP{\mathbf{P}}
\def\bR{\mathbf{R}}
\def\bS{\mathbf{S}}
\def\bT{\mathbf{T}}
\def\bX{\mathbf{X}}

% Bold numbers
\def\b0{\mathbf{0}}

% Bold greek
\bmdefine{\bmu}{\bm{\mu}}
\def\bphi{\bm{\phi}}
\def\bvarphi{\bm{\varphi}}

% Bold red sentence
\def\boldred#1{{\color{red}\textbf{#1}}}
\def\defword#1{{\color{orange}\textbf{#1}}}

% Caligraphic letters
\def\A{\mathcal{A}}
\def\B{\mathcal{B}}
\def\C{\mathcal{C}}
\def\D{\mathcal{D}}
\def\E{\mathcal{E}}
\def\F{\mathcal{F}}
\def\G{\mathcal{G}}
\def\H{\mathcal{H}}
\def\I{\mathcal{I}}
\def\L{\mathcal{L}}
\def\M{\mathcal{M}}
\def\N{\mathcal{N}}
\def\P{\mathcal{P}}
\def\R{\mathcal{R}}
\def\S{\mathcal{S}}
\def\T{\mathcal{T}}
\def\U{\mathcal{U}}
\def\V{\mathcal{V}}

% tt font for code
\def\code#1{{\tt #1}}

% i.e., e.g.
\def\eg{\emph{e.g.}}
\def\ie{\emph{i.e.}}


% Operators and special symbols
\def\nbOne{{\mathchoice {\rm 1\mskip-4mu l} {\rm 1\mskip-4mu l}
{\rm 1\mskip-4.5mu l} {\rm 1\mskip-5mu l}}}
\def\cov{\ensuremath{\mathsf{cov}}}
\def\Var{\ensuremath{\mathsf{Var}\ }}
\def\Im{\textrm{Im}\;}
\def\Re{\textrm{Re}\;}
\def\det{\ensuremath{\mathsf{det}}}
\def\diag{\ensuremath{\mathsf{diag}}}
\def\nullspace{\ensuremath{\mathsf{null}}}
\def\nullity{\ensuremath{\mathsf{nullity}}}
\def\rank{\ensuremath{\mathsf{rank}}}
\def\range{\ensuremath{\mathsf{range}}}
\def\sgn{\ensuremath{\mathsf{sgn}}}
\def\Span{\ensuremath{\mathsf{span}}}
\def\tr{\ensuremath{\mathsf{tr}}}
\def\imply{$\Rightarrow$}
\def\restrictTo#1#2{\left.#1\right|_{#2}}
\newcommand{\parallelsum}{\mathbin{\!/\mkern-5mu/\!}}
\def\dsum{\mathop{\displaystyle \sum }}%
\def\dind#1#2{_{\substack{#1\\ #2}}}

\DeclareMathOperator{\GL}{GL}
\DeclareMathOperator{\Rel}{Re}
\def\Nt#1{\left|\!\left|\!\left|#1\right|\!\right|\!\right|}
\newcommand{\tripbar}{|\! |\! |}



% The beamer bullet (in base colour)
\def\bbullet{\leavevmode\usebeamertemplate{itemize item}\ }

% Theorems and the like
\newtheorem{proposition}[theorem]{Proposition}
\newtheorem{property}[theorem]{Property}
\newtheorem{importantproperty}[theorem]{Property}
\newtheorem{importanttheorem}[theorem]{Theorem}
%\newtheorem{lemma}[theorem]{Lemma}
%\newtheorem{corollary}[theorem]{Corollary}
\newtheorem{remark}[theorem]{Remark}
\setbeamertemplate{theorems}[numbered]
%\setbeamertemplate{theorems}[ams style]

%
%\usecolortheme{orchid}
%\usecolortheme{orchid}

\def\red{\color[rgb]{1,0,0}}
\def\blue{\color[rgb]{0,0,1}}
\def\green{\color[rgb]{0,1,0}}


% Get rid of navigation stuff
\setbeamertemplate{navigation symbols}{}

% Set footline/header line
\setbeamertemplate{footline}
{%
\quad p. \insertpagenumber \quad--\quad \insertsection\vskip2pt
}
% \setbeamertemplate{headline}
% {%
% \quad\insertsection\hfill p. \insertpagenumber\quad\mbox{}\vskip2pt
% }


\makeatletter
\newlength\beamerleftmargin
\setlength\beamerleftmargin{\Gm@lmargin}
\makeatother

% Colours for special pages
\def\extraContent{yellow!20}


%%%%%%%%%%%%%%%%%
\usepackage{tikz}
\usetikzlibrary{shapes,arrows}
\usetikzlibrary{positioning}
\usetikzlibrary{shapes.symbols,shapes.callouts,patterns}
\usetikzlibrary{calc,fit}
\usetikzlibrary{backgrounds}
\usetikzlibrary{decorations.pathmorphing,fit,petri}
\usetikzlibrary{automata}
\usetikzlibrary{fadings}
\usetikzlibrary{patterns,hobby}

\usetikzlibrary{backgrounds,fit,petri}


\usepackage{pgfplots}
\pgfplotsset{compat=1.6}
\pgfplotsset{ticks=none}

\usetikzlibrary{decorations.markings}
\usetikzlibrary{arrows.meta}
\tikzset{>=stealth}

% For tikz
\usetikzlibrary{shapes,arrows}
\usetikzlibrary{positioning}
\tikzstyle{cloud} = [draw, ellipse,fill=red!20, node distance=0.87cm,
minimum height=2em]
\tikzstyle{line} = [draw, -latex']


%%% For max frame images
\newenvironment{changemargin}[2]{%
\begin{list}{}{%
\setlength{\topsep}{0pt}%
\setlength{\leftmargin}{#1}%
\setlength{\rightmargin}{#2}%
\setlength{\listparindent}{\parindent}%
\setlength{\itemindent}{\parindent}%
\setlength{\parsep}{\parskip}%
}%
\item[]}{\end{list}}


% Make one image take up the entire slide content area in beamer,.:
% centered/centred full-screen image, with title:
% This uses the whole screen except for the 1cm border around it
% all. 128x96mm
\newcommand{\titledFrameImage}[2]{
\begin{frame}{#1}
%\begin{changemargin}{-1cm}{-1cm}
\begin{center}
\includegraphics[width=108mm,height=\textheight,keepaspectratio]{#2}
\end{center}
%\end{changemargin}
\end{frame}
}

% Make one image take up the entire slide content area in beamer.:
% centered/centred full-screen image, no title:
% This uses the whole screen except for the 1cm border around it
% all. 128x96mm
\newcommand{\plainFrameImage}[1]{
\begin{frame}[plain]
%\begin{changemargin}{-1cm}{-1cm}
\begin{center}
\includegraphics[width=108mm,height=76mm,keepaspectratio]{#1}
\end{center}
%\end{changemargin}
\end{frame}
}

% Make one image take up the entire slide area, including borders, in beamer.:
% centered/centred full-screen image, no title:
% This uses the entire whole screen
\newcommand{\maxFrameImage}[1]{
\begin{frame}[plain]
\begin{changemargin}{-1cm}{-1cm}
\begin{center}
\includegraphics[width=\paperwidth,height=\paperheight,keepaspectratio]
{#1}
\end{center}
\end{changemargin}
\end{frame}
}

% This uses the entire whole screen (to include in frame)
\newcommand{\maxFrameImageNoFrame}[1]{
\begin{changemargin}{-1cm}{-1cm}
\begin{center}
\includegraphics[width=\paperwidth,height=0.99\paperheight,keepaspectratio]
{#1}
\end{center}
\end{changemargin}
}

% Make one image take up the entire slide area, including borders, in beamer.:
% centered/centred full-screen image, no title:
% This uses the entire whole screen
\newcommand{\maxFrameImageColor}[2]{
\begin{frame}[plain]
\setbeamercolor{normal text}{bg=#2!20}
\begin{changemargin}{-1cm}{-1cm}
\begin{center}
\includegraphics[width=\paperwidth,height=\paperheight,keepaspectratio]
{#1}
\end{center}
\end{changemargin}
\end{frame}
}


\usepackage{tikz}
\usetikzlibrary{patterns,hobby}
\usepackage{pgfplots}
\pgfplotsset{compat=1.6}
\pgfplotsset{ticks=none}

\usetikzlibrary{backgrounds}
\usetikzlibrary{decorations.markings}
\usetikzlibrary{arrows.meta}
\tikzset{>=stealth}

\tikzset{
  clockwise arrows/.style={
    postaction={
      decorate,
      decoration={
        markings,
        mark=between positions 0.1 and 0.9 step 40pt with {\arrow{>}},
   }}}}


   %%%%%%%%%%%
% To have links to parts in the outline
\makeatletter
\AtBeginPart{%
  \addtocontents{toc}{\protect\beamer@partintoc{\the\c@part}{\beamer@partnameshort}{\the\c@page}}%
}
%% number, shortname, page.
\providecommand\beamer@partintoc[3]{%
  \ifnum\c@tocdepth=-1\relax
    % requesting onlyparts.
    \makebox[6em]{Part #1:} \textcolor{green!30!blue}{\hyperlink{#2}{#2}}
    \par
  \fi
}
\define@key{beamertoc}{onlyparts}[]{%
  \c@tocdepth=-1\relax
}
\makeatother%

\newcommand{\nameofthepart}{}
\newcommand{\nupart}[1]%
    {   \part{#1}%
        \renewcommand{\nameofthepart}{#1}%
        {
          \setbeamercolor{background canvas}{bg=orange!50}
          \begin{frame}{#1}%\partpage 
          \hypertarget{\nameofthepart}{}\tableofcontents%
          \end{frame}
        }
    }



\usecolortheme{orchid}

%% Listings
\usepackage{listings}
\definecolor{mygreen}{rgb}{0,0.6,0}
\definecolor{mygray}{rgb}{0.5,0.5,0.5}
\definecolor{mymauve}{rgb}{0.58,0,0.82}
\definecolor{mygold}{rgb}{1,0.843,0}
\definecolor{myblue}{rgb}{0.537,0.812,0.941}

\definecolor{lgreen}{rgb}{0.6,0.9,.6}
\definecolor{lred}{rgb}{1,0.5,.5}

\lstloadlanguages{R}
\lstset{ %
  language=R,
  backgroundcolor=\color{black!05},   % choose the background color
  basicstyle=\footnotesize\ttfamily,        % size of fonts used for the code
  breaklines=true,                 % automatic line breaking only at whitespace
  captionpos=b,                    % sets the caption-position to bottom
  commentstyle=\color{mygreen},    % comment style
  escapeinside={\%*}{*)},          % if you want to add LaTeX within your code
  keywordstyle=\color{red},       % keyword style
  stringstyle=\color{mygold},     % string literal style
  keepspaces=true,
  columns=fullflexible,
  tabsize=4,
}
% Could also do (in lstset)
% basicstyle==\fontfamily{pcr}\footnotesize
\lstdefinelanguage{Renhanced}%
  {keywords={abbreviate,abline,abs,acos,acosh,action,add1,add,%
      aggregate,alias,Alias,alist,all,anova,any,aov,aperm,append,apply,%
      approx,approxfun,apropos,Arg,args,array,arrows,as,asin,asinh,%
      atan,atan2,atanh,attach,attr,attributes,autoload,autoloader,ave,%
      axis,backsolve,barplot,basename,besselI,besselJ,besselK,besselY,%
      beta,binomial,body,box,boxplot,break,browser,bug,builtins,bxp,by,%
      c,C,call,Call,case,cat,category,cbind,ceiling,character,char,%
      charmatch,check,chol,chol2inv,choose,chull,class,close,cm,codes,%
      coef,coefficients,co,col,colnames,colors,colours,commandArgs,%
      comment,complete,complex,conflicts,Conj,contents,contour,%
      contrasts,contr,control,helmert,contrib,convolve,cooks,coords,%
      distance,coplot,cor,cos,cosh,count,fields,cov,covratio,wt,CRAN,%
      create,crossprod,cummax,cummin,cumprod,cumsum,curve,cut,cycle,D,%
      data,dataentry,date,dbeta,dbinom,dcauchy,dchisq,de,debug,%
      debugger,Defunct,default,delay,delete,deltat,demo,de,density,%
      deparse,dependencies,Deprecated,deriv,description,detach,%
      dev2bitmap,dev,cur,deviance,off,prev,,dexp,df,dfbetas,dffits,%
      dgamma,dgeom,dget,dhyper,diag,diff,digamma,dim,dimnames,dir,%
      dirname,dlnorm,dlogis,dnbinom,dnchisq,dnorm,do,dotplot,double,%
      download,dpois,dput,drop,drop1,dsignrank,dt,dummy,dump,dunif,%
      duplicated,dweibull,dwilcox,dyn,edit,eff,effects,eigen,else,%
      emacs,end,environment,env,erase,eval,equal,evalq,example,exists,%
      exit,exp,expand,expression,External,extract,extractAIC,factor,%
      fail,family,fft,file,filled,find,fitted,fivenum,fix,floor,for,%
      For,formals,format,formatC,formula,Fortran,forwardsolve,frame,%
      frequency,ftable,ftable2table,function,gamma,Gamma,gammaCody,%
      gaussian,gc,gcinfo,gctorture,get,getenv,geterrmessage,getOption,%
      getwd,gl,glm,globalenv,gnome,GNOME,graphics,gray,grep,grey,grid,%
      gsub,hasTsp,hat,heat,help,hist,home,hsv,httpclient,I,identify,if,%
      ifelse,Im,image,\%in\%,index,influence,measures,inherits,install,%
      installed,integer,interaction,interactive,Internal,intersect,%
      inverse,invisible,IQR,is,jitter,kappa,kronecker,labels,lapply,%
      layout,lbeta,lchoose,lcm,legend,length,levels,lgamma,library,%
      licence,license,lines,list,lm,load,local,locator,log,log10,log1p,%
      log2,logical,loglin,lower,lowess,ls,lsfit,lsf,ls,machine,Machine,%
      mad,mahalanobis,make,link,margin,match,Math,matlines,mat,matplot,%
      matpoints,matrix,max,mean,median,memory,menu,merge,methods,min,%
      missing,Mod,mode,model,response,mosaicplot,mtext,mvfft,na,nan,%
      names,omit,nargs,nchar,ncol,NCOL,new,next,NextMethod,nextn,%
      nlevels,nlm,noquote,NotYetImplemented,NotYetUsed,nrow,NROW,null,%
      numeric,\%o\%,objects,offset,old,on,Ops,optim,optimise,optimize,%
      options,or,order,ordered,outer,package,packages,page,pairlist,%
      pairs,palette,panel,par,parent,parse,paste,path,pbeta,pbinom,%
      pcauchy,pchisq,pentagamma,persp,pexp,pf,pgamma,pgeom,phyper,pico,%
      pictex,piechart,Platform,plnorm,plogis,plot,pmatch,pmax,pmin,%
      pnbinom,pnchisq,pnorm,points,poisson,poly,polygon,polyroot,pos,%
      postscript,power,ppoints,ppois,predict,preplot,pretty,Primitive,%
      print,prmatrix,proc,prod,profile,proj,prompt,prop,provide,%
      psignrank,ps,pt,ptukey,punif,pweibull,pwilcox,q,qbeta,qbinom,%
      qcauchy,qchisq,qexp,qf,qgamma,qgeom,qhyper,qlnorm,qlogis,qnbinom,%
      qnchisq,qnorm,qpois,qqline,qqnorm,qqplot,qr,Q,qty,qy,qsignrank,%
      qt,qtukey,quantile,quasi,quit,qunif,quote,qweibull,qwilcox,%
      rainbow,range,rank,rbeta,rbind,rbinom,rcauchy,rchisq,Re,read,csv,%
      csv2,fwf,readline,socket,real,Recall,rect,reformulate,regexpr,%
      relevel,remove,rep,repeat,replace,replications,report,require,%
      resid,residuals,restart,return,rev,rexp,rf,rgamma,rgb,rgeom,R,%
      rhyper,rle,rlnorm,rlogis,rm,rnbinom,RNGkind,rnorm,round,row,%
      rownames,rowsum,rpois,rsignrank,rstandard,rstudent,rt,rug,runif,%
      rweibull,rwilcox,sample,sapply,save,scale,scan,scan,screen,sd,se,%
      search,searchpaths,segments,seq,sequence,setdiff,setequal,set,%
      setwd,show,sign,signif,sin,single,sinh,sink,solve,sort,source,%
      spline,splinefun,split,sqrt,stars,start,stat,stem,step,stop,%
      storage,strstrheight,stripplot,strsplit,structure,strwidth,sub,%
      subset,substitute,substr,substring,sum,summary,sunflowerplot,svd,%
      sweep,switch,symbol,symbols,symnum,sys,status,system,t,table,%
      tabulate,tan,tanh,tapply,tempfile,terms,terrain,tetragamma,text,%
      time,title,topo,trace,traceback,transform,tri,trigamma,trunc,try,%
      ts,tsp,typeof,unclass,undebug,undoc,union,unique,uniroot,unix,%
      unlink,unlist,unname,untrace,update,upper,url,UseMethod,var,%
      variable,vector,Version,vi,warning,warnings,weighted,weights,%
      which,while,window,write,\%x\%,x11,X11,xedit,xemacs,xinch,xor,%
      xpdrows,xy,xyinch,yinch,zapsmall,zip},%
   otherkeywords={!,!=,~,$,*,\%,\&,\%/\%,\%*\%,\%\%,<-,<<-,_,/},%
   alsoother={._$},%
   sensitive,%
   morecomment=[l]\#,%
   morestring=[d]",%
   morestring=[d]'% 2001 Robert Denham
  }%

%%%%%%% 
%% Definitions in yellow boxes
\usepackage{etoolbox}
\setbeamercolor{block title}{use=structure,fg=structure.fg,bg=structure.fg!40!bg}
\setbeamercolor{block body}{parent=normal text,use=block title,bg=block title.bg!20!bg}

\BeforeBeginEnvironment{definition}{%
	\setbeamercolor{block title}{fg=black,bg=yellow!20!white}
	\setbeamercolor{block body}{fg=black, bg=yellow!05!white}
}
\AfterEndEnvironment{definition}{
	\setbeamercolor{block title}{use=structure,fg=structure.fg,bg=structure.fg!20!bg}
	\setbeamercolor{block body}{parent=normal text,use=block title,bg=block title.bg!50!bg, fg=black}
}
\BeforeBeginEnvironment{importanttheorem}{%
	\setbeamercolor{block title}{fg=black,bg=red!20!white}
	\setbeamercolor{block body}{fg=black, bg=red!05!white}
}
\AfterEndEnvironment{importanttheorem}{
	\setbeamercolor{block title}{use=structure,fg=structure.fg,bg=structure.fg!20!bg}
	\setbeamercolor{block body}{parent=normal text,use=block title,bg=block title.bg!50!bg, fg=black}
}
\BeforeBeginEnvironment{importantproperty}{%
	\setbeamercolor{block title}{fg=black,bg=red!50!white}
	\setbeamercolor{block body}{fg=black, bg=red!30!white}
}
\AfterEndEnvironment{importantproperty}{
	\setbeamercolor{block title}{use=structure,fg=structure.fg,bg=structure.fg!20!bg}
	\setbeamercolor{block body}{parent=normal text,use=block title,bg=block title.bg!50!bg, fg=black}
}


% Beginning of a section
\AtBeginSection[]{
	{
		\setbeamercolor{background canvas}{bg=yellow!10}
		\begin{frame}[noframenumbering,plain]
			\framesubtitle{\nameofthepart Chapter \insertromanpartnumber \ -- \iteminsert{\insertpart}}
			\tableofcontents[
				currentsection,
				sectionstyle=show/shaded,
				subsectionstyle=show/show/hide]
		\end{frame}
	\addtocounter{page}{-1}
	%\addtocounter{framenumber}{-1} 
	}
}

% Beginning of a section
\AtBeginSubsection[]{
	{
		\setbeamercolor{background canvas}{bg=green!10}
		\begin{frame}[noframenumbering,plain]
				\framesubtitle{\nameofthepart Chapter \insertromanpartnumber \ -- \iteminsert{\insertpart}}
				\tableofcontents[
					currentsection,
					sectionstyle=show/show,
					currentsubsection,
					subsectionstyle=show/shaded/hide]
			\end{frame}
		\addtocounter{page}{-1}
	}
}


% To have theorems and everything derived prefixed by the slide set
% number
\renewcommand{\thetheorem}{5.\arabic{theorem}}



\title{Norms and Matrix Norms}
\author{Julien Arino}
\date{Fall 2023}

%%%%%%%%%%%%%%%%%%%%%
%%%%%%%%%%%%%%%%%%%%%
%%%%%%%%%%%%%%%%%%%%%
%%%%%%%%%%%%%%%%%%%%%
%%%%%%%%%%%%%%%%%%%%%
%%%%%%%%%%%%%%%%%%%%%
%%%%%%%%%%%%%%%%%%%%%
%%%%%%%%%%%%%%%%%%%%%
\begin{document}

% The title page
\begin{frame}[noframenumbering,plain]
  \begin{tikzpicture}[remember picture,overlay]
    \node[above right,inner sep=0pt] at (current page.south west)
    {
        \includegraphics[width=\paperwidth]{title-page-picture.png}
    };
\end{tikzpicture}
	\titlepage
\end{frame}
\addtocounter{page}{-1}
  
  
\begin{frame}{Outline}
	  \tableofcontents[hideallsubsections]
\end{frame}
\addtocounter{page}{-1}


%%%%%%%%%%%%%%%%%%%%%%%%%%%
%%%%%%%%%%%%%%%%%%%%%%%%%%%
%%%%%%%%%%%%%%%%%%%%%%%%%%%
%%%%%%%%%%%%%%%%%%%%%%%%%%%
\section{Vector norms}
\label{sec:vector_norms}


\begin{frame}
\begin{definition}[Norm]
Let $V$ be a vector space over a field $\IF$. A function $\| \cdot \|: V\to \IR_{+}$ is a \defword{norm} if for all $\bx,\by\in V$ and for all $c\in \IF$
\begin{enumerate}
    \item $\|\bx\| \geq 0$ \hfill[\defword{Nonnegativity}]
    \item $\|\bx\|=0 \iff\bx=\b0$ \hfill[\defword{Positivity}]
    \item $\|c\bx\| =|c|\ \|\bx\|$ \hfill[\defword{Homogeneity}]
    \item $\|\bx+\by \|\leq \|\bx\| + \|\by\|$ \hfill[\defword{Triangle Inequality}]
\end{enumerate}
\end{definition}
\vfill
\begin{remark}
If we have 1, 3, and 4 but not 2, then we have a \defword{seminorm}
\end{remark}
\end{frame}


\begin{frame}
\begin{definition}[Inner product]
Let $V$ be a vector space over $\IF$. A function $\langle\cdot,\cdot\rangle:  V \times V\to \IF$ is an \defword{inner product} if for all $\bx,\by,\bz\in V$ and all $c\in \IF$
\begin{enumerate}
    \item $\langle\bx,\bx\rangle \geq 0$
    \item $\langle\bx,\bx\rangle = 0 \iff \bx=0$
    \item $\langle\bx+\by,\bz\rangle=\langle\bx,\bz\rangle+\langle\by,\bz\rangle$
    \item $\langle c\bx,\by\rangle= c\langle\bx,\by\rangle$
    \item $\langle\bx,\by\rangle=\overline{\langle\by,\bx\rangle}$
\end{enumerate}
\end{definition}
\vfill
\begin{theorem}[Cauchy-Schwartz]
Let $\langle\cdot,\cdot\rangle$ be an inner product on a vector space $V$ over $\IF$, then 
\[
|\langle \bx,\by\rangle|^2
\leq \langle\bx,\bx\rangle\langle\by,\by\rangle
\]
\end{theorem}
\end{frame}


\begin{frame}
\begin{corollary}
If $\langle \cdot,\cdot \rangle$ is an inner product on a real or complex vector space $V$, then $\|\cdot\| : V\to \IR_{+}$ defined by $\| \bx \| = \langle \bx,\bx\rangle^{1/2}$ is a norm on $V$
\end{corollary}
\vfill
\begin{remark}
If $\langle \cdot,\cdot\rangle$ is a semi-inner product, then the resulting $\| \bx \| = \langle \bx, \bx\rangle^{1/2}$ is a seminorm
\end{remark}
\end{frame}



\begin{frame}
\begin{theorem}
 Consider the norm $\| \cdot \|$. Then $\| \cdot \|$ is derived from an inner product if and only if it satisfies the parallelogram identity
 \[\dfrac{1}{2}\left( \|\bx+\by\|^2 +\|\bx-\by\|^2\right)= \|\bx\|^2+\|\by\|^2\]
\end{theorem}

\begin{theorem}
If $\|\cdot\|$ is a nom on $\IC^n$ and a matrix $T\in \M_n$ which is non-singular. Then 
\[
    \|\bx\|_T= \| T\bx \|
\]
is also a norm on $\IC^n$
\end{theorem}
\end{frame}


%%%%%%%%%%%%%%%%%%%%%%%%%%%
%%%%%%%%%%%%%%%%%%%%%%%%%%%
%%%%%%%%%%%%%%%%%%%%%%%%%%%
%%%%%%%%%%%%%%%%%%%%%%%%%%%
\section{Analytic properties of norms}
\label{sec:analytic_properties_norms}


\begin{frame}
\begin{definition}
Let $V$ be a vector space over $\IF=\IR$ or $\IC$. Take a norm $\|\cdot\|$ on $V$. The sequence $\{\bx^{(k)}\}$ of vectors in $V$ converges to $\bx\in V$ with respect to the norm $\|\cdot\|$ if and only if $\| \bx^{(k)}-\bx\|\to 0$ as $k \to \infty$
\end{definition}
\vfill
We write $\lim\limits_{k \to \infty} \bx^{(k)}=\bx$ with respect to $\|\cdot\|$ or 
\[
    \bx^{(k)} \xrightarrow{\| \cdot \|} \bx 
\]
\end{frame}


\begin{frame}
\begin{theorem}
Every (vector) norm in $\IC^n$ is uniformly continuous
\end{theorem}
\vfill
\begin{corollary}
Let $\| \cdot \|_{\alpha}$ and $\| \cdot \|_{\beta}$ be any two norms on a finite-dimensional vector space $V$. Then there exist $C_m, C_r>0$ such that 
\[
    C_m\| \bx\|_{\alpha} \leq 
    \| \bx \|_{\beta}\leq 
    C_r\| \bx\|_{\alpha}, \forall \bx\in V
\]
\end{corollary}
\end{frame}


\begin{frame}
\begin{corollary}
 Let $\|\cdot\|_{\alpha}$ and $\|\cdot\|_{\beta}$ norms on a finite-dimensional vector space $V$ over $\IR$ or $\IC$, $\{\bx^{(k)}\}$ a given sequence in $V$, then 
 \[
    \bx^{(k)} \xrightarrow{\| \cdot \|_{\alpha}} \bx 
    \iff 
    \bx^{(k)} \xrightarrow{\| \cdot \|_{\beta}}\bx 
 \]
\end{corollary}
\end{frame}

\begin{frame}
\begin{definition}[Equivalent norms]
Two norms are \defword{equivalent}\index{equivalent} if whenever a sequence $\{\bx^{(k)}\}$ converges to $\bx$ with respect to one of the norm, it converges to $\bx$ in the other norm
\end{definition}
\vfill
\begin{theorem}
In finite-dimensional vector spaces, all norm are equivalent
\end{theorem}
\end{frame}



\begin{frame}
\begin{definition}[Dual norm]
Let $f$ be a pre-norm on $V=\IR^n$ or $\IC^n$. The function 
\[f_d=(\by) \max\limits_{f(\bx)=1} \Rel \by^* \bx\]
is the \defword{dual norm} of $f$
\end{definition}
\vfill
\begin{remark}
The dual norm is well defined. $\Rel\by^*\bx$ is a continuous function for all $\by\in V$ fixed. The set $\{f(\bx)=1\}$ is compact 
\end{remark}
\vfill
Equivalent definition for dual norm: $f^D(\by)= \max\limits_{f(\bx)=1} | \by^*\bx\mid$
\end{frame}


\begin{frame}
\begin{lemma}[Extension of Cauchy-Schwartz]
Let $f$ be a prenorm on $V=\IR^n$ or $\IC^n$ for all $\bx,\by\in V$. Then 
\begin{align*}
    & |\by^*\bx\mid \leq f(\bx)f^D(\by)\\
    & | \by^*\bx | \leq f^D(\bx)f(\bx)
\end{align*}
\end{lemma}
\vfill
\begin{remark}
\begin{itemize}
    \item The dual norm of a pre-norm is a norm
    \item The only norm that equals its dual norm is the Euclidean norm 
\end{itemize}
\end{remark}
\end{frame}



\begin{frame}
\begin{theorem}
Let $\|\cdot\|$ be a norm on $\IC^n$ or $\IR^n$, and $\|\cdot\|^D$ its dual, $c>0$ given. Then for all $\bx\in V$, $\|\bx\|= c\|\bx\|^d \iff \|\cdot\|= \sqrt{c}\|\cdot\|^d$.
In particular, $\|\cdot\|= \|\cdot\|^2 \iff \|\cdot\|= \|\cdot\|_2$
\end{theorem}
\vfill
\begin{definition}
Let $x\in \IF^n$. Denote $| x|= [|x_i|]$ ($|\cdot |$ entry-wise), and write that $|x|\leq |y |$ if $|x_i|\leq |y_i|$ for all $i= 1, \dots, n$. Assume $\|\cdot\|$ is 
\begin{enumerate}
    \item monotone if $|\bx|\leq |\by| \implies \|\bx\| \leq \|\by\|$ for all $\bx,\by$
    \item absloute if $\left \| |\bx| \right\| $ for all $\bx\in V$
\end{enumerate}
\end{definition}
\end{frame}



\begin{frame}
\begin{theorem}
Let $\|\cdot\|$ be a norm on $\IF^n$. Then
\begin{enumerate}
    \item If $\|\cdot\|$ is absolute, then 
    \[
        \|\by\|^D= \max\limits_{\bx\neq 0}= \dfrac{|\by|^T |\bx|}{\|\bx\|}
    \]
    for all $\by\in V$
    \item If $\|\cdot\|$ absolute, then $\|\cdot\|^D$ is absolute and monotone
    \item $\|\cdot\|$ absolute if and only if $\|\cdot\|$
\end{enumerate}
\end{theorem}
\end{frame}

%%%%%%%%%%%%%%%%%%%%%%%%%%%
%%%%%%%%%%%%%%%%%%%%%%%%%%%
%%%%%%%%%%%%%%%%%%%%%%%%%%%
%%%%%%%%%%%%%%%%%%%%%%%%%%%
\section{Matrix Norms}
\label{sec:matrix_norms}


\begin{frame}
\begin{definition}[Matrix norm]
Let $\Nt{\cdot} $ be a function from $\M_n\to \IR$. $\Nt{\cdot}$ is a \defword{matrix norm} if for all $A,B\in \M_n$ and $c\in\IC$, it satisfies the following
\begin{enumerate}
    \item $\Nt{A} \geq 0$ \hfill[\defword{nonnegativity}]
    \item $\Nt{A}  =0 \iff A=0$ \hfill[\defword{positivity}]
    \item $\Nt{cA} = |c|\ \Nt{A}$ \hfill[\defword{homogeneity}]
    \item $\Nt{A+B} \leq \Nt{A}  +\tripbar B \tripbar$ \hfill[\defword{triangle inequality}]
    \item $\tripbar AB \tripbar \leq \tripbar A\tripbar \tripbar B \tripbar$
    \hfill[\defword{submultiplicativity}]
\end{enumerate}
\end{definition}
\vfill
\begin{remark}
    As with vector norms, if property 2 does not hold, $\Nt{\cdot}$ is a \defword{matrix semi-norm}
\end{remark}
\end{frame} 



\begin{frame}

\begin{remark}
 $\tripbar A^2 \tripbar = \tripbar A A\tripbar\leq \tripbar A \tripbar^2 $ [for any matrix norm].\\
 
If $A^2= A$, then 
 \[\tripbar A^2 \tripbar = \tripbar A \tripbar\leq \tripbar A\tripbar^2 \implies \tripbar A \tripbar \geq 1.\]
 
 In particular, $\tripbar I \tripbar \geq 1$ for any matrix norm. \\
 
 Assume that $A$ is invertible, then $AA^{-1}=I$, thus 
 \begin{align}
     &\tripbar I \tripbar = \tripbar AA^{-1} \tripbar \leq \tripbar A \tripbar \tripbar A^{-1}\tripbar\\
     &\tripbar A^{-1} \tripbar\geq \dfrac{\tripbar I \tripbar }{\tripbar A\tripbar}
 \end{align}
\end{remark}
\end{frame}




\begin{frame}
\begin{definition}[Induced matrix norm]\label{def:induced_matrix_norm}
Let $\|\cdot\|$ be a norm on $\IC^n$. Define $\Nt{\cdot}$ on $\M_n(\IC)$ by 
\[
    \tripbar A\tripbar  = \max\limits_{\|\bx\|=1} \|A\bx\|
\]
Then $\Nt{\cdot}$ is the \defword{matrix norm induced} by $\|\cdot\|$
\end{definition}
\vfill
\begin{theorem}
The function $\Nt{\cdot}$ defined in Definition~\ref{def:induced_matrix_norm} has the following properties
\begin{enumerate}
    \item $\tripbar\II\tripbar =1$
    \item $\| A\by \| \leq \tripbar A \tripbar \|\by\|$ for all $A\in \M_n(\IC)$ and all $\by \in \IC^n$
    \item $\Nt{\cdot}$ is a matrix norm on $\M_n(\IC)$.
    \item $\tripbar A \tripbar = \max\limits_{\|\bx\|= \|\by\|^D} |\by^* A\bx|$
\end{enumerate}
\end{theorem}
\end{frame}


\begin{frame}
\begin{definition}[Induced norm/Operator norm]
$\Nt{\cdot}$ defined from $\|\cdot\|$ by any of the previous methods is the matrix norm induced by $\|\cdot\|$. It is also called the \defword{operator norm}
\end{definition}
\vfill
\begin{definition}[Unital norm]
A norm such that $\tripbar \mathbb{I} \tripbar=1$ is \defword{unital}
\end{definition}
\vfill
\begin{remark}
 Every induced matrix norm is unital. Every induced norm is a matrix norm
\end{remark}
\end{frame}


\begin{frame}
\begin{proposition}
For all $U,\, V$ unitary matrices, we have $\tripbar UAV\tripbar_2= \tripbar A\tripbar_2$
\end{proposition}
\vfill
\begin{theorem}
Let $\Nt{\cdot}$ be a matrix norm in $\M_n$ and let $S\in \M_n$ be nonsingular. 
Then for all $A\in \M_n$, $\tripbar A\tripbar_S= \tripbar SAS^{-1}\tripbar$ is a matrix norm. 
Furthermore, if $\Nt{\cdot}$ on $\IC^n$, then $\|\bx\|_S= \|S\bx\|$ induces $\Nt{\cdot}_S$ on $\M_n$
\end{theorem}
\vfill
\begin{theorem}
Let $\Nt{\cdot}$ be a matrix norm on $\M_n$, $A\in \M_n$ and $\lambda\in \sigma(A)$. 
Then 
\begin{enumerate}
    \item $|\lambda |\leq \rho(A) \leq \tripbar A\tripbar$
    \item If $A$ is nonsingular, then 
    \[
        \rho(A)\geq |\lambda |\geq \dfrac{1}{\tripbar A^{-1} \tripbar}
    \]
\end{enumerate}
\end{theorem}
\end{frame}



\begin{frame}
\begin{lemma}
Let $A\in \M_n$. If there exists a norm $\Nt{\cdot}$ on $\M_n$ such that $\tripbar A\tripbar< 1$, then $\lim\limits_{k \to \infty}A^k=0$ entry-wise
\end{lemma}
\vfill
\begin{remark}
 When $\tripbar A\tripbar <1$ for some norm, we say that $A$ is \defword{convergent}
\end{remark}
\vfill
\begin{theorem}
Let $A\in \M_n$, then 
\[
    \lim\limits_{k \to \infty} A^k=0 \iff \rho(A)<1
\]
\end{theorem}
\end{frame}



\begin{frame}
\begin{theorem}[Gelfand Formula] 
    Let $\Nt{\cdot}$ be a matrix norm on $\M_n$, let $A\in \M_n$. Then 
    \[
        \rho(A)= \lim\limits_{k \to \infty} \tripbar A^k \tripbar^{1/k}
    \]
\end{theorem}
\vfill
\begin{theorem}
 Let $R$ be the radius of convergence of the (scalar) power series $\sum\limits_{k=0}^{\infty} a_kz^k$ and $A\in \M_n$. 
 Then the matrix power series $\sum\limits_{k=1}^{\infty} a_k A^k$ converges if $\rho(A)<R$
\end{theorem}
\end{frame}


\begin{frame}
\begin{remark}
 The convergence condition for the matrix power series is ``there exists a matrix norm $\Nt{\cdot}$ such that $\tripbar A \tripbar <R$''
\end{remark}
\vfill
\begin{corollary}
  Let $A\in \M_n$ be nonsingular, if there $\Nt{\cdot}$ matrix norm such that $\tripbar \mathbb{I}-A\tripbar \leq 1$
\end{corollary}
\vfill
\begin{corollary}
  Let $A\in \M_n$ is such that $|a_{ii}| > \sum\limits_{j \neq i} |a_{ij}|$ for all $i = 1, \dots, n$. Then $A$ is invertible
\end{corollary}
\end{frame}



%%%%%%%%%%%%%%%%%%%%%%%%%%%
%%%%%%%%%%%%%%%%%%%%%%%%%%%
%%%%%%%%%%%%%%%%%%%%%%%%%%%
%%%%%%%%%%%%%%%%%%%%%%%%%%%
\section{Matrix norms and Singular values}
\label{sec:matrixNorms_SVD}

\begin{frame}
Let $V = \M_{mn}(\IC)$ with Frobenius inner product 
\[
    \langle A, \, B \rangle_{F}= \tr(B^*A)
\]
The norm derived from the Frobenius inner product is 
\[\|A\|_2= (\tr(A^*A))^{1/2}\]
is the $\ell$-2 norm (or Frobenius norm)
\end{frame}

\begin{frame}
The spectral norm $\Nt{\cdot}$ defined on $\M_n$ by 
\[\tripbar A\tripbar_2= \sigma_1(A),\]
where $\sigma_1(A)$ is the largest singular value of $A$ is induced by the $\ell$-2 norm on $\IC^n$. \\
Inded, from the singular value decomposition theorem, let 
\[A= V \Sigma W^*\]
be a singular value decomposition of $A$, where $V, \, W$ unitary, $\Sigma= \sigma(\sigma_1, \dots, \sigma_n)$ and $\sigma_1 \geq \dots \geq \sigma_n\geq 0$ are the non-increasingly ordered singular values of $A$
\end{frame}


\begin{frame}
From unitary invariance and monotonicity of the Euclidean norm, we say that 
\begin{align*}
\max \|Ax\|_1&= \max_{\|x\|_1} \|V\Sigma W^*\|_2\\
&= \max_{\|x\|_2} \|\Sigma W^* x\|_2\\
&= \max_{\|Wy\|_2=1}\|\Sigma y\|_2\\
&= \max_{\|y\|_2}\|\Sigma y\|_2\\
& \leq \max_{\|y\|_2}\|\sigma_1 y\|_2\\
&= \sigma_1 \max_{\| y \|_2} \|y\|_2\\
&= \sigma_1
\end{align*}
Since $\| \Sigma y\|_2= \sigma_1$ for $y= e_1$, 
\[\max_{\|x\|_2=1} \|Ax\|_2= \sigma_1(A)\]
\end{frame}


\begin{frame}
We could have used 
\begin{align*}
\max_{\|x\|_2=1}=\|Ax\|_2^2&= \max_{\|x\|_2=1}x^*A^*AX\\
&= \lambda_{max}(A^*A)\\
&= \sigma_1(A)
\end{align*}


\begin{remark}
	For all $U,\, V$ unitary $\M_{n}$ matrices, for all $A\in \M_n$, $\tripbar UAV\tripbar_2= \tripbar A\tripbar_2$
\end{remark}
\end{frame}



%%%%%%%%%%%%%%%%%%
%%%%%%%%%%%%%%%%%%
%%%%%%%%%%%%%%%%%%
%%%%%%%%%%%%%%%%%%
\begin{frame}[allowframebreaks]
    \frametitle{References}
    \bibliographystyle{amsalpha}
    \bibliography{MATH-4370-7370-lecture-notes.bib}
\end{frame}
    

\end{document}